\documentclass[twocolumn]{aastex631}
\begin{document}
\title{The reionization ionizing-photon-budget shortfall for star-forming galaxies is not robust to systematics at z\~{}9 (jwsthigh systematic corner)}
\author{NebulaMind Lab (autonomous pipeline)}
\affiliation{NebulaMind Lab, \url{https://lab.nebulamind.net}}
\begin{abstract}
Reionization ionizing-photon-budget reconciliation at z\~{}9: star-forming galaxies require f\_esc=0.180 (+0.170/-0.087) to close the budget (Madau-Dickinson SFRD, log xi\_ion=25.5+/-0.15, clumping C=2-5, JWST-SFRD tail), versus indirect-proxy-inferred f\_esc=0.062 (+0.110/-0.039) from LzLCS O32/beta calibrations. Median delta(required-inferred)=+0.103 dex-frac (16-84\%: -0.017 to +0.276); 81\% of the systematic MC shows a shortfall. Conclusion: the budget CLOSES within the systematic, and the sign holds under both O32 and beta calibrations. The result is bounded by the xi\_ion x clumping x proxy-calibration systematic, not by statistics. Generated autonomously from public data (jwst) via the NebulaMind Lab runner: a bounded, reproducible, descriptive study.
\end{abstract}
\section{Introduction}
Recent studies have highlighted a potential crisis in the reionization photon budget, suggesting that current observations may not account for the necessary ionizing photons to drive cosmic reionization [Muñoz2024]. This discrepancy has sparked interest in reconciling the photon budget using various approaches, including excursion set models [Park2022] and assessments of galaxy contributions [Duncan2015]. However, these efforts have yet to fully resolve the issue. Building on this work, we aim to investigate the reionization photon budget at z\~{}9 by comparing required ionizing photon production with indirect proxy-inferred values.
\section{Data and method}
To address this question, we adopt a literature-anchored budget calculation approach that does not rely on new survey catalog data. Instead, we use the Madau \& Dickinson (2014) analytic fitting function for the cosmic star formation rate density (SFRD). We also utilize published calibrations for the ionizing photon production efficiency (xi\_ion) and the escape fraction (f\_esc) from the Lyman-continuum galaxy sample (LzLCS) [Chisholm+22, Flury+22; Simmonds+24]. Our method involves calculating the required f\_esc to close the reionization photon budget based on these parameters.
\section{Result}
Our analysis yields a single quantitative result: star-forming galaxies at z\~{}9 require an escape fraction of f\_esc = 0.180 (+0.170/-0.087) to reconcile the ionizing photon budget, assuming the Madau-Dickinson SFRD, log xi\_ion=25.5±0.15, and a clumping factor C between 2-5. In contrast, indirect proxy-inferred values from LzLCS O32/beta calibrations suggest f\_esc = 0.062 (+0.110/-0.039). The median difference between the required and inferred escape fractions is +0.103 dex (16-84\% range: -0.017 to +0.276), with 81\% of systematic Monte Carlo simulations showing a shortfall in ionizing photons.
\begin{figure}\centering\includegraphics[width=\columnwidth]{result.png}\caption{The reionization ionizing-photon-budget shortfall for star-forming galaxies is not robust to systematics at z\~{}9 (jwsthigh systematic corner)}\end{figure}
\section{Caveats}
It is essential to acknowledge the limitations of our approach, which relies on automated, single-selection, uncalibrated measurements from published literature. The accuracy of our result depends heavily on the assumptions and calibrations used in previous studies, such as the Madau \& Dickinson (2014) SFRD fitting function and LzLCS O32/beta proxy calibrations. Additionally, our analysis does not account for potential systematic errors or uncertainties in these underlying measurements, which could impact the validity of our findings. Furthermore, the reliance on a single method and set of parameters may overlook alternative explanations or contributions to the reionization photon budget.
\section*{References}
\footnotesize
[Muoz2024] Muñoz et al. (2024). Reionization after JWST: a photon budget crisis?. bibcode 2024MNRAS.535L..37M \par [Park2022] Park et al. (2022). Calibrating excursion set reionization models to approximately conserve ionizing photons. bibcode 2022MNRAS.517..192P \par [Duncan2015] Duncan et al. (2015). Powering reionization: assessing the galaxy ionizing photon budget at z < 10. bibcode 2015MNRAS.451.2030D \par [Davies2021] Davies et al. (2021). The Predicament of Absorption-dominated Reionization: Increased Demands on Ionizing Sources. bibcode 2021ApJ...918L..35D \par [Madau2017] Madau (2017). Cosmic Reionization after Planck and before JWST: An Analytic Approach. bibcode 2017ApJ...851...50M

\end{document}
